%
% ---------------------------------------------------------------
% Copyright (C) 2012-2018 Gang Li
% ---------------------------------------------------------------
%
% This work is the default powerdot-tuliplab style test file and may be
% distributed and/or modified under the conditions of the LaTeX Project Public
% License, either version 1.3 of this license or (at your option) any later
% version. The latest version of this license is in
% http://www.latex-project.org/lppl.txt and version 1.3 or later is part of all
% distributions of LaTeX version 2003/12/01 or later.
%
% This work has the LPPL maintenance status "maintained".
%
% This Current Maintainer of this work is Gang Li.
%
%

\documentclass[
 size=14pt,
 paper=smartboard,  %a4paper, smartboard, screen
 mode=present, 		%present, handout, print
 display=slides, 	% slidesnotes, notes, slides
 style=tuliplab,  	% TULIP Lab style
 pauseslide,
 fleqn,leqno]{powerdot}


\usepackage{amssymb}
\usepackage{amsmath}
\usepackage{rotating}
\usepackage{graphicx}
\usepackage{boxedminipage}
\usepackage{media9}
\usepackage{rotate}
\usepackage{calc}
\usepackage[absolute]{textpos}
\usepackage{psfrag,overpic}
\usepackage{fouriernc}
\usepackage{pstricks,pst-node,pst-text,pst-3d,pst-grad}
\usepackage{moreverb,epsfig,subfigure}
\usepackage{pstricks}
\usepackage{pstricks-add}
\usepackage{pst-text}
\usepackage{pst-node, pst-tree}
\usepackage{booktabs}
\usepackage{etex}
\usepackage{breqn}
\usepackage{multirow}
\usepackage{gitinfo2}
\usepackage{xcolor}

\usepackage{todonotes}
% \usepackage{pst-rel-points}

\usepackage{listings}
\lstset{frameround=fttt,
frame=trBL,
stringstyle=\ttfamily,
backgroundcolor=\color{yellow!20},
basicstyle=\footnotesize\ttfamily}
\lstnewenvironment{code}{
\lstset{frame=single,escapeinside=`',
backgroundcolor=\color{yellow!20},
basicstyle=\footnotesize\ttfamily}
}{}


\usepackage{hyperref}
\hypersetup{ % TODO: PDF meta Data
  pdftitle={Presentation Title},
  pdfauthor={Gang Li},
  pdfpagemode={FullScreen},
  pdfborder={0 0 0}
}


% \usepackage{auto-pst-pdf}
% package to show source code

\definecolor{LightGray}{rgb}{0.9,0.9,0.9}
\newlength{\pixel}\setlength\pixel{0.000714285714\slidewidth}
\setlength{\TPHorizModule}{\slidewidth}
\setlength{\TPVertModule}{\slideheight}
\newcommand\highlight[1]{\fbox{#1}}
\newcommand\icite[1]{{\footnotesize [#1]}}

\newcommand\twotonebox[2]{\fcolorbox{pdcolor2}{pdcolor2}
{#1\vphantom{#2}}\fcolorbox{pdcolor2}{white}{#2\vphantom{#1}}}
\newcommand\twotoneboxo[2]{\fcolorbox{pdcolor2}{pdcolor2}
{#1}\fcolorbox{pdcolor2}{white}{#2}}
\newcommand\vpspace[1]{\vphantom{\vspace{#1}}}
\newcommand\hpspace[1]{\hphantom{\hspace{#1}}}
\newcommand\COMMENT[1]{}

\newcommand\placepos[3]{\hbox to\z@{\kern#1
        \raisebox{-#2}[\z@][\z@]{#3}\hss}\ignorespaces}

\renewcommand{\baselinestretch}{1.2}


\newcommand{\draftnote}[3]{
	\todo[author=#2,color=#1!30,size=\footnotesize]{\textsf{#3}}	}
% TODO: add yourself here:
%
\newcommand{\gangli}[1]{\draftnote{blue}{GLi:}{#1}}
\newcommand{\shaoni}[1]{\draftnote{green}{sn:}{#1}}
\newcommand{\gliMarker}
	{\todo[author=GLi,size=\tiny,inline,color=blue!40]
	{Gang Li has worked up to here.}}
\newcommand{\snMarker}
	{\todo[author=Sn,size=\tiny,inline,color=green!40]
	{Shaoni has worked up to here.}}

%%%%%%%%%%%%%%%%%%%%%%%%%%%%%%%%%%%%%%%%%%%%%%%%%%%%%%%%%%%%%%%%%%%%%%%%
% title
% TODO: Customize to your Own Title, Name, Address
%
\title{Team for Universal Learning and Intelligent Processing}
\author{
Gang Li
\\
\\School of Information Technology
\\Deakin University, Australia
}
\date{\gitCommitterDate}


% Customize the setting of slides
\pdsetup{
% TODO: Customize the left footer, and right footer
rf=\href{http://www.tulip.org.au}{
Last Changed by: \textsc{\gitCommitterName}\ \gitVtagn-\gitAbbrevHash\ (\gitAuthorDate)
},
cf={Recent Research \@ TULIP Lab},
}


\begin{document}

\maketitle

%\begin{slide}{Overview}
%\tableofcontents[content=sections]
%\end{slide}


%%==========================================================================================
%%
\begin{slide}[toc=,bm=]{Overview}
\tableofcontents[content=currentsection,type=1]
\end{slide}
%%
%%==========================================================================================


\section{TULIP Lab}


%%==========================================================================================
%%
\begin{slide}{TULIP Team}
\begin{center}
\twotonebox{\rotatebox{90}{Defn}}{\parbox{.86\textwidth}
{Outlying Aspects Mining aims to
identify the outstanding features of the query object.
\begin{itemize}
\item A teacher may be interested in the \textcolor{orange}{characteristics} that
make \textcolor{orange}{one student} \textcolor{orange}{distinctive} from others.
\item NBA coaches would prefer to
find out the strengths and weaknesses of the player (a query object).
\end{itemize}
}}

\end{center}
\bigskip
\begin{center}
\begin{tabular}{c| c c c c }
\toprule
Player & \texttt{3PT\%}  & \texttt{FTA} & \texttt{FT\%} & \texttt{To} \\
\midrule
$P_1$
&  {$65$} &  {$4$} &  {$33$} &  {$8$} \\
$P_2$
&  {$78$} &  {$1$}&  {$65$}&  {$5$} \\
$P_3$
&  {$58$} &  {$6$} &  {$46$} &  {$3$} \\
$P_4$
&  {$68$} &  {$1.2$}&  {$85$}&  {$6.2$} \\
$P_5$
&  {$58$} &  {$6.2$} &  {$36$} &  {$3.4$}\\
\bottomrule
\end{tabular}
\end{center}

\end{slide}
%%
%%==========================================================================================


%%==========================================================================================
%%
\begin{slide}{TULIP Lab}
\begin{center}
\begin{tabular}{c| c c c c }
\toprule
%\centering
Player & \texttt{3PT\%}  & \texttt{FTA} & \texttt{FT\%} & \texttt{To} \\
\midrule
$P_1$
&  {$65$} &  {$4$} &  {$33$} &  {$8$} \\
$P_2$
&  {$78$} &  {$1$}&  {$65$}&  {$5$} \\
$P_3$
&  {$58$} &  {$6$} &  {$46$} &  {$3$} \\
$P_4$
&  {$68$} &  {$1.2$}&  {$85$}&  {$6.2$} \\
$P_5$
&  {$58$} &  {$6.2$} &  {$36$} &  {$3.4$}\\
\bottomrule
\end{tabular}
\end{center}

\bigskip

\twocolumn[
\savevalue{lfrheight}=4.6cm,
\savevalue{lfrprop}={
linestyle=solid,framearc=.2,linewidth=1pt},
rfrheight=\usevalue{lfrheight},
rfrprop=\usevalue{lfrprop}
]{
Outlying Aspects Mining
\begin{itemize}
\item
\smallskip
Explain the distinctive \textcolor{orange}{aspects} of the query object.
\smallskip
\item
\smallskip
The query object may (or may not) be an outlier.
\end{itemize}
}{
Outlier Detection
\begin{itemize}
\item
\smallskip
Find out \textcolor{orange}{all} unusual
\textcolor{orange}{objects} in the whole dataset.
\smallskip
\item
\smallskip
\textcolor{orange}{No} explanation on how they are different.
\end{itemize}
}

\end{slide}
%%
%%==========================================================================================


%%==========================================================================================
%%
\begin{slide}{Lab Organization}
\twotonebox {\rotatebox{90}{Defn}}{\parbox{.88\textwidth}
{
{\textcolor{orange}{Group outlying aspects mining} aims to
identify the outstanding features of the group of query object.
\begin{itemize}
\item
Doctors desire to identify the merits \& demerits between
\textcolor{orange}{a group of cancer patients} and normal people.
\item
NBA coaches are passionate about exploring the obvious advantages \&
disadvantages of \textcolor{orange}{the team}.
\end{itemize}
}
}}

\vspace{1.5cm}

\twocolumn{
\begin{figure}
  \centering
  \selectcolormodel{rgb}
  \missingfigure{Testing.}
  %\includegraphics[width=0.6\textwidth]{figures//demical.eps}\\
  \caption{Medical}\label{fig:demical}
\end{figure}
}{
\begin{figure}
  \centering
  \selectcolormodel{rgb}
  \missingfigure{Testing.}
  %\includegraphics[width=0.6\textwidth]{figures//NBA_team.eps}\\
  \caption{NBA-Team}\label{fig:timg}
\end{figure}
}

\end{slide}
%%
%%==========================================================================================


\section{Research at TULIP}


%%==========================================================================================
%%
\begin{slide}[toc=,bm=]{Research at TULIP}
\begin{itemize}
    \item 
    Current Research Themes @ TULIP Lab
    \begin{itemize}
        \item 
        \textcolor{orange}{Behavior Informatics}
            \begin{itemize}
                \item 
                Periodic Behavior Mining
                
                \item 
                Behavior Prediction
            \end{itemize}
            
        \item 
        \textcolor{orange}{Information Abuse Prevention}
            \begin{itemize}
                \item 
                Data Compliance Checking
                
                \item 
                Privacy Preservation in Data Mining  
            \end{itemize}
        \item 
        \textcolor{orange}{Business Intelligence Applications}
            \begin{itemize}
                \item 
                Recommender System
                
                \item 
                Tourism/Hospitality Management        
            \end{itemize}
    \end{itemize}
\end{itemize}
\end{slide}
%%
%%==========================================================================================


%%==========================================================================================
%%
\begin{slide}{Theme 1. Behavior Informatics}
%Theme one -
\begin{itemize}
\item
People share \textcolor{green}{news, interests} 
and \textcolor{green}{ideas} in OSNs,

\begin{itemize}
\item
though these platforms also spread 
\textcolor{red}{email malware, rumours, gossips}
\textcolor{red}{and malicious links}, and also leak our 
\textcolor{red}{privacy}.
\end{itemize}
\end{itemize}

\end{slide}
%%
%%==========================================================================================


%%==========================================================================================
%%
\begin{slide}[toc=,bm=]{The Popularity of Online Social Networks}

\begin{itemize}
\item 
The popularity of OSNs has greatly increased in recent years

\begin{itemize}
\item
A large amount of Internet users also use OSN services,

\item
Australia has up to 11 million Facebook users.
\end{itemize}

\end{itemize}
\end{slide}
%%
%%==========================================================================================


%%==========================================================================================
%%
\begin{slide}[toc=,bm=]{Application of Geo Photos}

\begin{itemize}
\item
There are many geotagged photos available. However :

\begin{itemize}
\item
Data is noisy or misleading.

\item
Photos are taken in transit rather than at the attractions.

\item
Photo is static media, whereas, travel behavior is dynamic.
\end{itemize}

\end{itemize}

\end{slide}
%%
%%==========================================================================================


%%==========================================================================================
%%
\begin{slide}[toc=,bm=]{Tourist Movement Analysis}

Tourist Traffic Flow in Hong Kong Metropolitan Area.

\end{slide}
%%
%%==========================================================================================


%%==========================================================================================
%%
\begin{slide}[toc=,bm=]{Tourist Movement Analysis}

Huy Quan Vu, Gang Li, Rob Law, Ben Haobin Yip. Exploring the travel behaviors of inbound tourists to Hong Kong using Geotagged photos. Tourism Management.


\end{slide}
%%
%%==========================================================================================


%%==========================================================================================
%%
\begin{slide}[toc=,bm=]{Tourist Movement Analysis}


\end{slide}
%%
%%==========================================================================================


%%==========================================================================================
%%
\begin{slide}[toc=,bm=]{Periodic Behavior Mining}

\begin{itemize}
\item 
Observed Check-in Data is usually a mixed events from different periodic behaviors

\begin{itemize}
\item
What are Periodic?

\item
What is the period?

\item
Predict the Behavior on a particular day.
\end{itemize}

\end{itemize}

\end{slide}
%%
%%==========================================================================================


%%==========================================================================================
%%
\begin{slide}[toc=,bm=]{Periodic Behavior Mining}

\begin{itemize}
\item
Observed Check-in Data is usually a mixed events from different periodic behaviors

    \begin{itemize}
    \item
    What are Periodic?

    \item
    What is the period?

    \item
    Predict the Behavior on a particular day.

    \end{itemize}
\end{itemize}

\end{slide}
%%
%%==========================================================================================


%%==========================================================================================
%%
\begin{slide}[toc=,bm=]{Periodic Behavior Mining Mobility Intentions}

\begin{itemize}
\item
CP Decomposition

\end{itemize}

\end{slide}
%%
%%==========================================================================================


%%==========================================================================================
%%
\begin{slide}[toc=,bm=]{Periodic Behavior Mining Applications}

\begin{itemize}
\item 
There are various applications of Periodic Behavior Mining

    \begin{itemize}
    \item
    What is the schedule of a person on a particular day/time?

    \item
    Where to find the person on a particular day/time?
        \begin{itemize}
        \item
        Police applications

        \item
        Scheduled ``Encountering'' for match making
        \end{itemize}
    \item
    Timely Recommendation Systems
        \begin{itemize}
        \item
        Now it is the time for you to stock rice, cooking oil, or ice cream!

        \end{itemize}

\end{itemize}

\end{itemize}

\end{slide}
%%
%%==========================================================================================


%%==========================================================================================
%%
\begin{slide}[toc=,bm=]{Data Science Application - Case 1: Face Age Recognition}


\end{slide}
%%
%%==========================================================================================


%%==========================================================================================
%%
\begin{slide}[toc=,bm=]{Data Science Application - Case 2: Speech-based Emotion Detection}


\end{slide}
%%
%%==========================================================================================


%%==========================================================================================
%%
\begin{slide}[toc=,bm=]{Data Science Application - Case 3: K-Complex Detection}


\end{slide}
%%
%%==========================================================================================


%%==========================================================================================
%%
\begin{slide}[toc=,bm=]{Data Science Application - Case 4: Trajectory Analysis}


\end{slide}
%%
%%==========================================================================================


%%==========================================================================================
%%
\begin{slide}[toc=,bm=]{Data Science Application - Case 5: Video Surveillance}


\end{slide}
%%
%%==========================================================================================


%%==========================================================================================
%%
\begin{slide}[toc=,bm=]{Market Segmentation Analysis}


\end{slide}
%%
%%==========================================================================================

%%==========================================================================================
%%
\begin{slide}{Theme 2. Information Abuse Prevention}
%Theme two -
\begin{itemize}
\item
Information Abuse:


\begin{itemize}
\item
The compromise of information 
for which the data stakeholders are not willing to disclose.


\item
It depends on the nature of the information, 
the involved internal or external users, 
and the ways in which the organization grants the access to 
or releases the information.

\end{itemize}
\end{itemize}

\end{slide}
%%
%%==========================================================================================


%%==========================================================================================
%%
\begin{slide}[toc=,bm=]{Who is an adversary?}

\begin{itemize}
\item
Every user is potentially an adversary

\begin{itemize}
\item
After data is released, we cannot prevent any user from performing any type of analysis on the released data

\item
Worst case scenario

\begin{itemize}
\item 
Must account for disclosure risk from any and all types of analyses
\end{itemize}

\end{itemize}
\end{itemize}

\end{slide}
%%
%%==========================================================================================


%%==========================================================================================
%%
\begin{slide}[toc=,bm=]{Information Privacy Applications - Case 1: Privacy Breach}

\end{slide}
%%
%%==========================================================================================


%%==========================================================================================
%%
\begin{slide}[toc=,bm=]{Information Privacy Applications - Case 2: Privacy Breach}


\end{slide}
%%
%%==========================================================================================


%%==========================================================================================
%%
\begin{slide}[toc=,bm=]{Information Privacy Applications - Case 3: Is O.S.N. a Secure Place to Show off?}


\end{slide}
%%
%%==========================================================================================

%%==========================================================================================
%%
\begin{slide}[toc=,bm=]{Privacy Preserving Related Reference}


\end{slide}
%%
%%==========================================================================================

%%==========================================================================================
%%
\begin{slide}[toc=,bm=]{Theme 3. Business Intelligence}

\begin{itemize}
\item 
Tourism Data Mining

\begin{itemize}
\item
Tourist Behavior Analysis

\item
Contrast Mining

\begin{itemize}
\item
Outlying Aspects Mining

\item
Competitiveness Discovery
\end{itemize}

\end{itemize}

\item
Recommendation System

\end{itemize}

\end{slide}
%%
%%==========================================================================================


%%==========================================================================================
%%
\begin{slide}[toc=,bm=]{Business Intelligence Applications - Case 1: Market Segmentation Analysis}


\end{slide}
%%
%%==========================================================================================


%%==========================================================================================
%%
\begin{slide}[toc=,bm=]{Business Intelligence Applications - Case 2: Negative Association Rule Mining}


\end{slide}
%%
%%==========================================================================================


%%==========================================================================================
%%
\begin{slide}[toc=,bm=]{Business Intelligence Applications - Case 3: Contrast Mining}


\end{slide}
%%
%%==========================================================================================


%%==========================================================================================
%%
\begin{slide}[toc=,bm=]{Business Intelligence Applications - Case 4: Heterogeneous Preferences}

\begin{itemize}
\item
User-wise biases

\begin{itemize}
\item
Users give lower ratings when in bad mood

\item
Different scales, e.g., 4/5 star is ``OK'' for critics

\item
User interface design

\end{itemize}

\item
Item-wise biases

\begin{itemize}
\item 
The same item presented differently, 
e.g., a movie in 3D, HD

\end{itemize}
\end{itemize}
\end{slide}
%%
%%==========================================================================================


%%==========================================================================================
%%
\begin{slide}[toc=,bm=]{Business Intelligence Applications - Case 5: PR-based Conditional Random Fields}


\end{slide}
%%
%%==========================================================================================

%%==========================================================================================
%%
\begin{slide}[toc=,bm=]{Our rankings}


\end{slide}
%%
%%==========================================================================================

%%==========================================================================================
%%
\begin{slide}[toc=,bm=]{Awards}


\end{slide}
%%
%%==========================================================================================


\section{Training at TULIP}


%%==========================================================================================
%%
\begin{slide}{Training at TULIP}

Framework of GOAM algorithm:

\bigskip

\begin{figure}
  \centering
  \selectcolormodel{rgb}
  \missingfigure{Testing a long text string.}
%  \includegraphics[width=0.55\textwidth]{figures//framework1.eps}\\
  \caption{Framework of GOAM Algorithm} \label{framework}
\end{figure}

\end{slide}
%%
%%==========================================================================================


%%==========================================================================================
%%
\begin{slide}[toc=,bm=]{Training at TULIP}



\end{slide}
%%
%%==========================================================================================


%%==========================================================================================
%%
\begin{slide}{Moving to Higher Stages}
\begin{itemize}
\item
\smallskip
Suppose $f_1$, $f_2$, $f_3$ are three features of $G_q$.

$f_1$: \{$x_1, x_2, x_3, x_4, x_5, x_2, x_3, x_4, x_1, x_2$\} \\

$f_2$: \{$y_2, y_2, y_1, y_2, y_3, y_3, y_5, y_4, y_4, y_2$\} \\

$f_3$: \{$z_1, z_4, z_2, z_4, z_5, z_3, z_1, z_2, z_4, z_2$\} \\
\end{itemize}

\begin{figure}[htbp]
    \centering
    \subfigure[$f_1$]{
        \selectcolormodel{rgb}
        \missingfigure[figwidth=5.5cm]{Test.}
        %\includegraphics[width=0.25\textwidth]{figures//frequency-distribution-feature1.eps}
        \label{fig:fre-dis-f1}
    }
    \subfigure[$f_2$]{
        \selectcolormodel{rgb}
        \missingfigure[figwidth=5.5cm]{Test.}
        \label{fig:fre-dis-f2}
    }
    \subfigure[$f_3$]{
        \selectcolormodel{rgb}
        \missingfigure[figwidth=5.5cm]{Test.}
        \label{fig:fre-dis-f3}
    }
    \caption{Histogram of $G_q$ on three features}
    \label{fig:fre-dis-each-feature}
\end{figure}

\end{slide}
%%
%%==========================================================================================


\section{Rules in TULIP}


%%==========================================================================================
%%
\begin{slide}{Regulations}

\begin{center}
\begin{itemize}

\item<1->
{TULIP Academy}

\begin{itemize}
\item
You must not be absent from the meeting without permission.

\item
Xi'an group members must attend the seminar in the Lab,
and NOT online.

\item
Collaborate with team members, rather than attacking each other.
\end{itemize}

\item<2->
{FLIP Study}

\begin{itemize}
\item
All FLIP students need to pass the corresponding prerequisite.

\item
All FLIP materials are confidential, and you should not circulate it outside.
\end{itemize}

\item<3->
{Research}

\begin{itemize}
\item
SmartGit, Bitbucket, Mendeley, update PPR in time.

\item
Actively discuss research work in the group meeting, and TULIP Seminar.

\item
Everyone should prepare at least one question for each TULIP seminar.
\end{itemize}
\end{itemize}
\end{center}

\end{slide}
%%
%%==========================================================================================


%%==========================================================================================
%
\begin{slide}[toc=,bm=]{Questions?}



\end{slide}
%%
%%==========================================================================================


%%==========================================================================================
% TODO: Contact Page
\begin{wideslide}[toc=,bm=]{Contact Information}
\centering
\vspace{\stretch{1}}
\twocolumn[
lcolwidth=0.35\linewidth,
rcolwidth=0.65\linewidth
]
{
% \centerline{\includegraphics[scale=.2]{tulip-logo.eps}}
}
{
\vspace{\stretch{1}}
Associate Professor \emph{Gang Li}\\
School of Information Technology\\
Deakin University, Geelong, Australia
\begin{description}
 \item[Email] \href{mailto:gangli@acm.org}
 {\textsc{\footnotesize{gangli@acm.org}}}

 \item[Lab] \href{http://www.tulip.org.au}
 {\textsc{\footnotesize{Team for Universal Learning and Intelligent Processing}}}
\end{description}
}
\vspace{\stretch{1}}
\end{wideslide}

\end{document}

\endinput
